\section{Background and Motivation}
\label{sec:background}

\subsection{Real-world agent work}

We use \emph{real-world agent work} to mean work in which agents help users
manage concrete tasks rather than only produce isolated model outputs. Such
work may involve local state, external repositories, remote terminals, GUI and
TUI controls, model calls, package releases, documents, and human decisions.
The product problem is not only how to run an agent. It is how to keep the
state of the work legible over time.

Long-lived work introduces a persistence problem. Users and agents need to ask:
what happened, what source observed it, what evidence supports it, what is
missing, and what action is safe now? A timeline view is a natural interface
for these questions, but only if the view makes its own assumptions inspectable.

\subsection{The global-clock temptation}

Distributed systems theory has long shown that time and ordering are not
trivial in distributed computation. Lamport's happened-before relation
separates causal order from physical clock order \cite{lamport1978time};
vector-clock style systems make concurrency explicit \cite{mattern1989virtual}.
Yet many product surfaces still sort mixed-source facts primarily by wall-clock
timestamps and present the result as a single history. That is often sufficient
for casual logs, but it is not sufficient for recovery, audit, or cooperation
among capable participants.

The problem is sharper in agent systems because the user may not have the time
or expertise to manually diagnose missing evidence. A local product must help
both the human and the agent understand what the view means.

\subsection{KFD foundation}

The KFD foundation used by Kungfu can be summarized as:

\begin{quote}
KFD-1: facts must not drift.\\
KFD-2: trust must start from facts.\\
KFD-3: cooperation must start from trusted value.
\end{quote}

KFD-4 applies this foundation to timelines:

\begin{quote}
KFD-4: timelines must declare their observer.
\end{quote}

Under this framing, a timeline view is trustworthy only when it is grounded in
non-drifting facts, exposes enough evidence for trust assessment, and helps
participants cooperate through transparent value rather than hidden pressure or
opaque authority \cite{kfd2026observer}.

\subsection{Why a causal object is needed}

A projection rule alone does not define the unit that can be accepted, moved,
verified, or rejected. Raw journal pages are physical append blocks; sessions
and runs are product vocabulary; individual frames are too small to provide a
portable closure boundary. Kungfu therefore introduces an \emph{Episode}: a
bounded causal segment containing the manifest, frame references, content
commitments, provenance, schemas, dependencies, and verification evidence
needed to inspect one unit of work \cite{kungfu2026episode}.

Episodes do not make the timeline absolute. They make its inputs explicit. A
view can name which Episodes it accepted, preserve their internal and declared
cross-Episode causality, and then apply observer policy only where the causal
facts leave an ordering choice.
